\subsection{\texorpdfstring{The spaces \(\mathcal{L}_{\mathfrak{S}}(\mathbf{E};\mathbf{F})\)}{The spaces \textbackslash mathcal\{L\}\_\{\textbackslash mathfrak\{S\}\}(\textbackslash mathbf\{E\};\textbackslash mathbf\{F\})}}

Let \(\mathbf{F}\) be a topological vector space, \(\mathbf{E}\) an arbitrary set, and \(\mathfrak{S}\) a family of subsets of \(\mathbf{E}\) . Consider the vector space \(\mathrm{F}^{\mathrm{E}}\) with the uniform structure of \(\mathfrak{S}\) -convergence (GT, X, § 1, No.~2). We know that this structure is compatible with the commutative group structure of \(\mathrm{F}^{\mathrm{E}}\) (GT, X, § 1, No.~4, cor. 2). The topology so deduced is called the \(\mathfrak{S}\) -topology. If \(\mathbf{X}\) is a subset of \(\mathrm{F}^{\mathrm{E}}\) , or more generally, a set with a mapping \(j: \mathrm{X} \to \mathrm{F}^{\mathrm{E}}\) , then the inverse image under \(j\) of the \(\mathfrak{S}\) -topology on \(\mathrm{F}^{\mathrm{E}}\) is called the \(\mathfrak{S}\) -topology on \(\mathbf{X}\) .

Remarks. --- 1) The \(\mathfrak{S}\) -topology is identical with the \(\mathfrak{S}'\) -topology, where \(\mathfrak{S}'\) denotes the bornology generated by \(\mathfrak{S}\) (III, p.~1).

\begin{enumerate}
\def\labelenumi{\arabic{enumi})}
\setcounter{enumi}{1}
\tightlist
\item
  Let \(\mathbf{M} \in \mathfrak{S}\) and let \(\mathbf{V}\) be a neighbourhood of 0 in \(\mathbf{F}\) ; let \(\mathrm{T}(\mathbf{M}, \mathbf{V})\) denote the set of all \(f \in \mathrm{F}^{\mathrm{E}}\) such that \(f(x) \in \mathbf{V}\) for every \(x \in \mathbf{M}\) . If \(\mathfrak{S}\) is stable under finite unions, the sets \(\mathrm{T}(\mathbf{M}, \mathbf{V})\) form a fundamental system of neighbourhoods of 0 for the \(\mathfrak{S}\) -topology of \(\mathrm{F}^{\mathrm{E}}\) .
\end{enumerate}

PROPOSITION 1. --- Let E be a set, S a family of subsets of E, F a topological vector space and H a vector subspace of \(F^{E}\) . In order that the S-topology be compatible with the vector space structure of H, it is necessary and sufficient that \(u(\mathbf{M})\) is bounded in F for every \(u \in H\) and every \(M \in S\) . If, moreover, F is locally convex, then the S-topology on H is locally convex.

On account of Remarks 1) and 2) above, we see that a necessary and sufficient condition for the \(\mathfrak{S}\) -topology to be compatible with the vector space structure of H is that the sets \(\mathrm{H} \cap \mathrm{T}(\mathbf{M}, \mathrm{V})\) are absorbent in H (I, p.~7, prop. 4); but this implies that for every \(u \in \mathrm{H}\) , every subset \(\mathbf{M} \in \mathfrak{S}\) , and every balanced neighbourhood V of 0 in F, there exists \(\lambda \neq 0\) such that \(u(M) \subset \lambda V\) ; that is to say (III, p.~2) that \(u(\mathbf{M})\) is bounded in F. Finally, the last assertion of the proposition follows from the fact that if V is convex, so is \(\mathrm{T}(\mathbf{M}, \mathrm{V})\) .

COROLLARY. --- Let E and F be two locally convex spaces, S a family of bounded subsets of E, and \(\mathcal{L}(E; F)\) the vector space of continuous linear mappings from E into F. Then the S-topology is compatible with the vector space structure of \(\mathcal{L}(E; F)\) and is locally convex.

It is enough to remark that if \(u\) is a continuous linear mapping from \(\mathbf{E}\) into \(\mathbf{F}\) and \(\mathbf{M}\) is a bounded subset of \(\mathbf{E}\) , then \(u(\mathbf{M})\) is bounded in \(\mathbf{F}\) (III, p.~4, cor. 1).

Given two locally convex vector spaces E and F, and a family S of bounded subsets of E, let \(\mathcal{L}_{\mathfrak{S}}(E;F)\) denote the locally convex space obtained by assigning the S-topology to \(\mathcal{L}(E;F)\) .

Examples. --- 1) If \(\mathfrak{S}\) is the set of all finite subsets of E, then the \(\mathfrak{S}\) -topology is the topology of simple convergence and the space \(\mathcal{L}_{\mathfrak{S}}(E; F)\) is also denoted by \(\mathcal{L}_s(E; F)\) . A bounded subset of \(\mathcal{L}_s(E; F)\) is called a simply bounded subset of \(\mathcal{L}(E; F)\) .

\begin{enumerate}
\def\labelenumi{\arabic{enumi})}
\setcounter{enumi}{1}
\item
  If \(\mathfrak{S}\) is the set of compact (resp. precompact, compact convex) subsets, then the \(\mathfrak{S}\) -topology is called the topology of compact (resp. precompact, compact convex) convergence and the space \(\mathcal{L}_{\mathfrak{S}}(\mathrm{E};\mathrm{F})\) is also denoted by \(\mathcal{L}_c(\mathrm{E};\mathrm{F})\) (resp. \(\mathcal{L}_{pc}(\mathrm{E};\mathrm{F}),\mathcal{L}_{cc}(\mathrm{E};\mathrm{F})\) ). (Cf. IV, p.~48, exerc. 7.)
\item
  If \(\mathfrak{S}\) is the set of all bounded subsets of E, we say that the \(\mathfrak{S}\) -topology is the topology of bounded convergence and the space \(\mathcal{L}_{\mathfrak{S}}(\mathrm{E};\mathrm{F})\) is denoted by \(\mathcal{L}_b(\mathrm{E};\mathrm{F})\) .
\item
  When \(\mathrm{F} = \mathbf{K}\) , the space \(\mathcal{L}(\mathrm{E};\mathrm{F})\) is the dual \(\mathrm{E}'\) of \(\mathrm{E}\) . We denote by \(\mathrm{E}_{\Xi}^{\prime}, \mathrm{E}_{s}^{\prime}\) etc. the space \(\mathcal{L}_{\Xi}(\mathrm{E};\mathbf{K}), \mathcal{L}_s(\mathrm{E};\mathbf{K})\) etc. The space \(\mathrm{E}_s^\prime\) (resp. \(\mathrm{E}_b^\prime\) ) is called the weak dual (resp. strong dual) of \(\mathrm{E}\) . A bounded subset of \(\mathrm{E}_s^\prime\) (resp. \(\mathrm{E}_b^\prime\) ) is said to be weakly (resp. strongly) bounded. We observe that the weak topology on \(\mathrm{E}'\) is none other than \(\sigma (\mathrm{E}',\mathrm{E})\) (II, p.~42).
\end{enumerate}

When \(\mathrm{E} = \mathrm{F}\) , we denote by \(\mathcal{L}(\mathrm{E}), \mathcal{L}_{\mathfrak{S}}(\mathrm{E})\) etc. the space \(\mathcal{L}(\mathrm{E};\mathrm{F}), \mathcal{L}_{\mathfrak{S}}(\mathrm{E};\mathrm{F})\) etc. Let \(p\) be a continuous semi-norm on \(\mathrm{F}\) and \(\mathbf{M}\) a bounded subset of \(\mathrm{E}\) . Let

\[
p _ {\mathrm{M}} (u) = \sup _ {x \in \mathrm{M}} p (u (x)).\tag{1}
\]

It is immediate that \(p_{M}\) is a semi-norm on \(\mathcal{L}(E;F)\) and that if \(\Gamma\) is a fundamental system of semi-norms on F, the family of semi-norms \(p_{M}\) , where p ranges over \(\Gamma\) and M ranges over a base for the bornology generated by S, is a fundamental system of semi-norms of \(\mathcal{L}_{\mathfrak{S}}(E;F)\) .

In particular, if E and F are semi-normed spaces, and if p (resp. q) denotes the semi-norm of E (resp. F), then the topology of bounded convergence on \(\mathcal{L}(\mathrm{E};\mathrm{F})\) is defined by the semi-norm

\[
r (u) = \sup _ {p (x) \leqslant 1} q (u (x))\tag{2}
\]

(cf.~GT, X, § 3, No.~2). When we consider \(\mathcal{L}_b(\mathrm{E};\mathrm{F})\) as a semi-normed space, we shall always, unless the contrary is expressly stated, mean the semi-norm (2). If F is a normed space, the semi-norm (2) is a norm.

Remarks. --- 3) Let A be a dense subset of the unit ball of E. In view of the continuity of u, we also have

\[
r (u) = \sup _ {x \in \mathrm{A}} q (u (x)).\tag{3}
\]

For example

\[
r (u) = \sup _ {p (x) <   1} q (u (x)).\tag{4}
\]

Since we have \(u(tx) = tu(x)\) for \(t \in \mathbf{R}\) , we also have,

\[
r (u) = \sup _ {p (x) = 1} q (u (x)) = \sup _ {p (x) \neq 0} \frac {q (u (x))}{p (x)}.\tag{5}
\]

whenever \(p \neq 0\) .

\begin{enumerate}
\def\labelenumi{\arabic{enumi})}
\setcounter{enumi}{3}
\tightlist
\item
  The formula (2) shows that the map \(u \mapsto r(u)\) is lower semi-continuous on \(\mathcal{L}_s(\mathrm{E};\mathrm{F})\) .
\end{enumerate}

PROPOSITION 2. --- Let E and F be two locally convex spaces and let S be a set of bounded subsets of E.

\begin{enumerate}
\def\labelenumi{\arabic{enumi})}
\item
  The \(\mathfrak{S}\) -topology on \(\mathcal{L}(\mathrm{E};\mathrm{F})\) is identical with the \(\mathfrak{S}\) -topology, where \(\mathfrak{S}\) denotes the smallest adapted, bornology (III, p.~3) on \(\mathrm{E}\) which contains \(\mathfrak{S}\) .
\item
  Suppose that \(\{0\}\) is not dense in \(\mathbf{F}\) and let \(\mathfrak{S}'\) be another set of bounded subsets of \(\mathbf{E}\) . Then the \(\mathfrak{S}'\) -topology is coarser than the \(\mathfrak{S}\) -topology if and only if \(\mathfrak{S}' \subset \tilde{\mathfrak{S}}\) .
\end{enumerate}

Let \(u \in \mathcal{L}(\mathrm{E};\mathrm{F})\) , \(\mathbf{M} \in \mathfrak{S}\) and let \(p\) be a continuous semi-norm on \(\mathrm{F}\) . Since \(p \circ u\) is a continuous semi-norm on \(\mathrm{E}\) , this is the same as saying that \(p \circ u\) is bounded above by 1 on \(\mathbf{M}\) or on the closed, convex balanced envelope \(\tilde{\mathbf{M}}\) of \(\mathbf{M}\) ; in other words, we have \(p_{\mathbf{M}} = p_{\tilde{\mathbf{M}}}\) . Moreover, it is clear that we have \(p_{\lambda \mathbf{M}} = \lambda p_{\mathbf{M}}\) for all \(\lambda > 0\) and \(p_{\mathbf{M} \cup \mathbf{M}'} = \sup(p_{\mathbf{M}}, p_{\mathbf{M}'})\) , from which the first assertion follows, since \(\tilde{\mathfrak{S}}\) has the set of homothetics of the closed convex balanced envelopes of finite unions of sets of \(\mathfrak{S}\) as a base.

We now prove the second assertion: first, if \(\mathbf{F}\) is the base field, it follows from the definition that the \(\tilde{\mathfrak{S}}\) -topology on \(\mathrm{E}' = \mathcal{L}(\mathrm{E};\mathrm{F})\) has as a fundamental system of neighbourhoods of 0, the set of polars of the sets of \(\tilde{\mathfrak{S}}\) . Let \(\mathbf{A}\) be a bounded subset of \(\mathbf{E}\) , whose polar \(\mathbf{A}^{\circ}\) is a neighbourhood of 0 for the \(\tilde{\mathfrak{S}}\) -topology; then there exists a closed convex balanced set \(\mathbf{B} \in \tilde{\mathfrak{S}}\) such that \(\mathbf{A}^{\circ} \supset \mathbf{B}^{\circ}\) , and so \(\mathbf{A} \subset \mathbf{B}^{\circ \circ}\) ; but by cor. 3 of II, p.~45, we have \(\mathbf{B}^{\circ \circ} = \mathbf{B}\) , and hence \(\mathbf{A} \subset \mathbf{B}\) and \(\mathbf{A} \in \tilde{\mathfrak{S}}\) . Therefore if \(\mathfrak{S}'\) is a set of bounded subsets of \(\mathbf{E}\) , the \(\mathfrak{S}'\) -topology is coarser than the \(\mathfrak{S}\) -topology on \(\mathrm{E}'\) if and only if \(\mathfrak{S}' \subset \tilde{\mathfrak{S}}\) . The general case follows immediately, since if \(y \in F\) is not in the closure of 0, we can verify that the mapping which makes \(f \in E'\) correspond to the mapping \(x \mapsto f(x)y\) is an isomorphism of the locally convex spaces \(\mathrm{E}'_{\mathfrak{S}}\) onto its image in \(\mathcal{L}_{\mathfrak{S}}(\mathrm{E};\mathrm{F})\) .

